% Part: computability
% Chapter: computability-theory
% Section: total

\documentclass[../../../include/open-logic-section]{subfiles}

\begin{document}

\olfileid{cmp}{thy}{tot}
\olsection{Totalitas Ora Bisa Diputusake}

Ayo kita gatekake conto liyane panganggone teorema $s$-$m$-$n$
kanggo nuduhake yen sawijining perkara ora komputabel. Tetepna
$\fn{Tot}$ minangka himpunan indeks fungsi komputabel total, yaiku
\[
\fn{Tot} = \Setabs{x}{\text{kanggo saben $y$, $\cfind{x}(y)\fdefined$}}.
\]

\begin{prop}
\ollabel{prop:total}
$\fn{Tot}$ ora komputabel.
\end{prop}

\begin{proof}
Kanggo ndeleng yen $\fn{Tot}$ ora komputabel, cukup tuduhna yen $K$
bisa direduksi menyang himpunan iku. Tetepna $h(x,y)$ kang ditetepake dening
\[
h(x,y) \simeq
\begin{cases}
0 & \text{yen $x \in K$} \\
\fundefined & \text{yen ora}
\end{cases}
\]
Gatekna yen $h(x,y)$ babar pisan ora gumantung marang $y$.
Ora angel ndeleng yen $h$~komputabel parsial: kanggo
input $x, y$, kita ngitung~$h$ kanthi luwih dhisik
niru fungsi~$\cfind{x}$ nganggo input~$x$; yen komputasi
iki mandheg, $h(x,y)$ ngasilake $0$ lan mandheg. Mula
$h(x,y)$ padha karo $\Zero(\umin{s}{T(x,x,s)})$, ing ngendi
$\Zero$ minangka fungsi konstan nol.

Kanthi nggunakake teorema $s$-$m$-$n$, ana fungsi
rekursif primitif~$k(x)$ saengga kanggo saben $x$ lan~$y$,
\[
\cfind{k(x)}(y) =
\begin{cases}
0 & \text{yen $x \in K$} \\
\fundefined & \text{yen ora}
\end{cases}
\]
Mula $\cfind{k(x)}$ total yen $x \in K$, lan ora
ditegesi ing input endi wae yen ora. Dadi,
$k$ minangka reduksi $K$ menyang~$\fn{Tot}$.
\end{proof}

\begin{digress}
Pranyata $\fn{Tot}$ malah ora bisa dienumerasi sacara
komputabel---kompleksitase luwih dhuwur ing
``hierarki aritmetika.'' Nanging kita ora bakal
ngrembug penguatan iki ing kene.
\end{digress}

\end{document}
